\section{Conventions, standing hypotheses and concordance}\label{S:conventions}

\subsection{Standing conventions}
The main text fixes four conventions in Section~2 and uses them throughout; they are restated here with the places where each one is load-bearing.

\begin{description}[leftmargin=1.6em,itemsep=2pt]
\item[(Fail)] An action outside $U(x)$ at a compatible state $x$ is failure. Used in the admissibility case of Theorem~4.1, in clause~(i) of the selector principle (Proposition~2.2, proof in \S\ref{S:selector}), and in the minimal construction of Proposition~4.5, whose state-dependent control set $U(S)=\{0,r(S)\}$ is a modelling primitive carrying this convention.
\item[(Hold)] Reviews occur at intervals of positive length $\Delta$, and the command selected at a review is held over the review interval. Theorem~4.1 and Proposition~4.2 are statements about this class; Theorem~4.6 and Theorem~4.8 name their own blind-window class; the coverage audit of Section~7 is restricted to the whole-window hold class.
\item[(Record)] The information state is the set of current states compatible with the record and the applied controls (main~2); the policy is a measurable nonanticipative function of the record. A disturbance realisation is a measurable selection $d(t)\in D(x(t))$ along the trajectory it generates (memoryless model); concatenations of selections along a glued trajectory are again selections.
\item[(Order)] Viability is ``exists policy, for all compatible states and admissible realisations'': the record is generated by the realisation, so conditioning the admissible realisations on the record would be circular.
\end{description}

\paragraph{Regularity.} $f$ is continuous, $U$ and $D$ are compact-valued with closed graph, and Carath\'eodory solutions exist for every admissible input pair on the horizons considered; $\cV=\bigcap_j\{q_j\ge0\}$ with $q_j\in C^1$. Under nonunique solutions every property ranges over every solution of the closed-loop pair. No universal sensor model is posited (main text, Section~8.3); each certificate names the observation structure it uses.

\paragraph{Adverse-path hypotheses.} Three certificates assume a realised original-system path: \hyp{C3} of Theorem~4.1, \hyp{D2} of Theorem~4.6, and \hyp{E2} of Theorem~4.10. In each case the hypothesis is on a genuine trajectory of the original game, not on a convexified inclusion, and \S\ref{S:phantom} shows by example that it does not follow from the closed graph or the lower semicontinuity of $D$.

\subsection{Concordance with the previous edition}
Table~\ref{tab:concordance} maps the results of this edition to the previous edition (v72) of the manuscript. Statements marked ``unchanged'' are mathematically identical up to notation; statements marked ``sharpened'' carry an added clause, a narrowed hypothesis or a corrected label; statements marked ``new'' have no v72 counterpart.

\begin{table}[htbp]
\centering\small
\begin{tabularx}{\textwidth}{@{}l>{\raggedright\arraybackslash}X>{\raggedright\arraybackslash}X@{}}
\toprule
v73 & content & v72 counterpart / status\\
\midrule
Def.~2.1 & robust epistemic kernel, hierarchy (main~1) & Def.~1, Section~2.3; unchanged\\
Prop.~2.2 & selector principle (four correspondences) & Prop.~1, Section~2.4; unchanged, proof \S\ref{S:selector}\\
Thm.~3.1 & finite-horizon soundness and completeness & Thm.~1, Section~3.1; unchanged\\
Thm.~3.2 & one-step completeness; obstruction tree & Thm.~2; unchanged\\
Ex.~3.3 & hidden-mode conflict & Ex.~1; unchanged\\
Thm.~4.1 & common-action obstruction \hyp{C1}--\hyp{C3} & Thm.~3 (H2.1--H2.3); unchanged\\
Prop.~4.2 & obstruction ladder & Prop.~2; label corrected to ``declared hold class''\\
Prop.~4.3 & sparse witness ($m{+}1$ states) & Prop.~3; unchanged, proof \S\ref{S:helly}\\
Rem.~4.4 & nonconvex counterexample, window rank & Section~3.4 remark; unchanged\\
Prop.~4.5 & minimal admissibility construction & Prop.~4; unchanged\\
Thm.~4.6 & delayed-information obstruction \hyp{D1}--\hyp{D3} & Thm.~4 (H3.1--H3.3); unchanged\\
Rem.~4.7 & threshold form $\sigma^*_\Pi$ & Rem.~2; unchanged\\
Thm.~4.8 & LP instantiation, (i)--(iv) & Thm.~5; unchanged, proof \S\ref{S:lp}\\
Ex.~4.9 & decaying threshold; hold-class gap & Section~3.3 examples; the two models now kept explicitly distinct\\
Thm.~4.10 & finite-time exit \hyp{E1}--\hyp{E2} & Thm.~6 (H5.1--H5.2); unchanged\\
Rem.~4.11 & comparison-function form & Rem.~3; unchanged\\
Prop.~4.12 & oracle-recourse obstruction $\Gamma_h$ & Prop.~5; support-function form, unchanged; control-affine counterexample \S\ref{S:recourse}\\
Ex.~4.13 & three-branch instance, $\tau^*=7/50$ & Section~3.7 instance; unchanged\\
Def.~5.1--Cor.~5.3 & exact certifier, fibre criterion, $\cY_{\mathrm{safe}}$ & Def.~2, Prop.~6, Cor.~1; unchanged\\
Prop.~5.4 & certainty-equivalence trap & Rem.~4; promoted to a proposition, statement unchanged\\
Thm.~6.1 & static-observation reduction & Thm.~7; sharpened (Nagumo test sufficient, not necessary)\\
Prop.~6.2 & two-phase decomposition & Prop.~8; unchanged\\
Prop.~6.3 & window no-go & Prop.~9; unchanged\\
Prop.~6.4 & monotonicity & Prop.~7; unchanged\\
Open Problem~6.5 & dynamic-observation completeness & Section~7 open problem; unchanged\\
Section~7 & resource model, coverage audit, two-patch table & Section~9; unchanged data, Table~1 now 48 cells\\
Table~3 & certificate $\leftrightarrow$ correspondence $\leftrightarrow$ remedy & Section~6.1 table; remedy column new\\
\S\ref{S:prob} & probabilistic counterpart (Thm.~8, Props.~10--11 of v72) & Section~10; moved here unchanged\\
\S\ref{S:landscape} & sufficiency landscape & Section~5 / S2; moved here unchanged\\
\S\ref{S:constructions} & constructions A.1--A.3 & S3; unchanged\\
\S\ref{S:verification} & Lean record; verification record V1--V12 & Section~3.9 + script; record consolidated\\
\bottomrule
\end{tabularx}
\caption{Concordance between the v73 and v72 editions.}
\label{tab:concordance}
\end{table}
